\pageTitre{INTRODUCTION}

Le Cameroun connaît depuis plusieurs années une forte croissance de la construction : immeubles
d'habitation, hôtels, bureaux et équipements publics se multiplient dans les grandes villes.
Chacun de ces bâtiments doit être équipé d'installations techniques fiables : distribution
d'eau, évacuation des eaux usées, climatisation, production de froid. Ces corps d'état
secondaires représentent une part importante du coût d'un ouvrage et conditionnent directement
le confort, l'hygiène et la sécurité de ses occupants. Ils créent aussi une demande soutenue en
techniciens qualifiés.

Apprenant en spécialité \textbf{\InfoSpecialite} au \InfoCFP, nous avons choisi d'effectuer
notre stage de fin de formation au sein de \InfoEntreprise{} (\InfoEntrepriseSigle) pour trois
raisons. D'abord, l'entreprise intervient à la fois en climatisation, en plomberie et en
électricité : c'est un cadre idéal pour comprendre comment ces lots s'articulent sur un même
chantier. Ensuite, elle exécute de grands projets, ce qui permet de découvrir l'organisation
d'un chantier d'envergure. Enfin, nous souhaitions consolider nos compétences en travail des
tubes, en assemblage et en essais d'étanchéité, qui sont communes au froid et à la plomberie.

Durant \InfoDuree, nous avons été affecté au chantier \InfoChantier, un immeuble de type
\InfoChantierNiveaux, où \InfoEntrepriseSigle{} assure notamment le lot plomberie sanitaire.
Le thème retenu pour ce rapport est donc la \textbf{plomberie sanitaire}, et plus précisément
la \emph{\MakeLowercase{\InfoSousTheme}}. Ce sujet est d'un intérêt certain : un réseau mal
dimensionné ou mal réalisé provoque des débits insuffisants aux étages, des coups de bélier,
des fuites et des surconsommations. Il se situe en outre à la frontière de notre spécialité,
puisque la distribution d'eau chaude, l'isolation des canalisations et la production d'eau
glacée font appel aux mêmes principes.

Pour rendre compte de cette expérience, le présent rapport s'articule autour de quatre parties.
La \textbf{première partie} présente l'entreprise d'accueil, son organisation et son
environnement. La \textbf{deuxième partie} décrit le déroulement du stage : poste occupé,
planning des activités, résultats et compétences acquises. La \textbf{troisième partie} traite
du thème technique : généralités sur la plomberie sanitaire, dimensionnement d'une colonne de
distribution, matériaux, mise en œuvre et liens avec le froid et la climatisation. La
\textbf{quatrième partie} est consacrée aux critiques et suggestions. Une conclusion générale
dresse le bilan du stage et ouvre sur nos perspectives professionnelles.
